\section{Statistical objective and the residual decomposition}

The target sample identifies its own treatment effect. Sources are used to
improve precision, while allowing an unknown set of arbitrarily small source
biases and a source count $K=K_n$ that may increase. We do not first require
consistent classification of every source. Parameter compatibility, rather
than source selection, is the inferential object. This follows the motivation
of searching under local invalidity \citep{guo2023}; it is not an application
of that paper's IV theorem. Target-specific borrowing, unknown source validity,
and increasing source counts are also connected to FACE, RIFL, and robust
fusion \citep{face2025,rifl2025,fusion2023}.

Condition on an independent training sample $\calT$. Fit a common prediction
$m^a$ and each source's merged weight $q_j^a$. Define
\[
 d=\E_t(m^1-m^0),\qquad r_a=\E_t\{Y(a)-m^a(X)\},\qquad
 \tau_t=d+r_1-r_0.
\]
Source $j$ evaluates $R_{ji}^a=\ind(A_{ji}=a)q_j^a(X_{ji})\{Y_{ji}-m^a(X_{ji})\}$
on fresh patients and sends its mean $\bar R_j^a$, unbiased mean variance
$\hat v_j^a=S_j^{2,a}/n_j$, and sample size. Its conditional mean is
$\theta_j^a=r_a+\beta_j^a$. Training, pilot, and evaluation are sample
\emph{roles} in this prototype, not a claim that the production RoCE
algorithm has changed from two-layer cross-fitting.

At least $g_a\ge1$ sources are valid in arm $a$; their identities may differ
between arms. Validity removes structural transport bias, not necessarily
fitting drift. Fix some truly valid subset $V_a$ of size $g_a$. We require
one of two certified forms of nuisance control:
\begin{align}
 &|\beta_j^a|\le b_a\quad(j\in V_a),\label{eq:h-uniform}\\
 &\left|g_a^{-1}\sum_{j\in V_a}\beta_j^a\right|\le B_{1a},\qquad
 K^{-1}\sum_{j\in V_a}(\beta_j^a)^2\le B_{2a}.
 \label{eq:h-average}
\end{align}
The second is a signed subset-mean and a squared-drift certificate. An
average over all sources does not establish it. All nuisance certificates,
including the target score drift, share a declared failure allowance
$\delta_B$. No minimum nonzero structural bias is assumed.

\section{Two complementary source constraints}

Write $\rho_a=g_a/K$. The source set is $C_a=J_a\cap F_a$ in the hybrid
procedure. Dispersion-only and Fourier-only procedures are prespecified
ablations with their own error allocations. Intersecting two separately
95\% procedures would not produce a 95\% hybrid.

\subsection{Dispersion and its direct bounded-score calibration}

Suppress $a$, let $\bar\theta=K^{-1}\sum_j\theta_j$, and define
\[
 D=K^{-1}\sum_j(\theta_j-\bar\theta)^2,\qquad
 \widehat D=K^{-1}\sum_j(\bar R_j-\bar R_\cdot)^2
       -\frac{K-1}{K^2}\sum_j\hat v_j.
\]
Independent patients within and between sites give $\E\widehat D=D$.
The statistic can be negative and is not truncated before inversion.
Subset geometry gives
\begin{equation}
 |\bar\theta-r|\le B_1+\sqrt{(K-g)D/g}.
 \label{eq:h-geometry}
\end{equation}
Under \eqref{eq:h-uniform}, take $B_1=b$.
Indeed, the difference between the full mean and a $g$-subset mean is
bounded by $\sqrt{(K-g)D/g}$ by Cauchy--Schwarz.

The earlier Cantelli construction remains a comparison: with certified
variance upper bounds and $\Var(\widehat D)\le LD+Q$, invert
$D-\widehat D\le\sqrt{(1-\alpha_D)(LD+Q)/\alpha_D}$.
Here the Cantelli constants use $L=4\max_jv_{Uj}/K$ and
$Q=2K^{-4}[(\sum_jv_{Uj})^2+\sum_j\{n_j(K-1)^2/(n_j-1)-1\}v_{Uj}^2]$.
The following bounded-score construction avoids individual variance
certificates altogether.

\begin{proposition}[Direct exponential dispersion bound]
Suppose patients are independent, identically distributed within each site,
$n_j\ge2$, and each site's score has a training-determined range of length
$R_j$. Put
\begin{align*}
 p_j&=R_j^2/(4n_j),&
 L&=4\max_jp_j/K,\\
 Q&=\frac{2}{K^4}\left[(\sum_jp_j)^2+
 \sum_j\left\{\frac{n_j}{n_j-1}(K-1)^2-1\right\}p_j^2\right],&
 c&=2\max_j\frac{(K-1)p_j}{K^2(n_j-1)}.
\end{align*}
For $x=\log(1/\alpha_D)$ and $y=\widehat D+cx$, define
\[
 U_D=\max\{0,y+xL+\sqrt{\max(0,2xLy+x^2L^2+2xQ)}\}.
\]
Then $\Prb(D\le U_D\mid\calT)\ge1-\alpha_D$. For $K=1$, set $D=U_D=0$.
\end{proposition}

\begin{proof}
Let $A=I/K-\boldsymbol1\boldsymbol1^T/K^2$ and form a patient matrix $B$.
Its within-site off-diagonal entries are $A_{jj}/\{n_j(n_j-1)\}$;
between sites they are $A_{jk}/(n_jn_k)$; its diagonal is zero.
For centered patient scores $X$, expansion gives
$\widehat D-D=2a^TX+X^TBX$, where $a_{ji}=(A\theta)_j/n_j$.
Each centered score is sub-Gaussian with proxy $s_j^2=R_j^2/4$ by
Hoeffding's lemma. Because the exponent is affine in each coordinate
separately, successive conditional MGF bounds replace independent $X_i$
by independent $N(0,s_j^2)$ variables without decreasing the MGF of
$-t(\widehat D-D)$.

Set $C=\operatorname{diag}(s)B\operatorname{diag}(s)$ and
$b=\operatorname{diag}(s)a$. On the site-constant subspace $C$ is equivalent
to $\operatorname{diag}(\sqrt p)A\operatorname{diag}(\sqrt p)\succeq0$.
On each site's orthogonal complement its eigenvalue is
$-A_{jj}p_j/(n_j-1)$. Also $b$ lies in the constant subspace,
$\operatorname{tr}(C)=0$, $2\operatorname{tr}(C^2)=Q$, and
$4\|b\|^2\le LD$. Gaussian integration and
$u-\log(1+u)\le u^2/\{2(1-ct)\}$ for the relevant eigenvalues give,
for $0<t<1/c$,
\[
 \log\E\exp\{-t(\widehat D-D)\}
 \le\frac{t^2(LD+Q)}{2(1-ct)}.
\]
Here the linear Gaussian term is at most $2t^2\|b\|^2$, since the
constant block is positive semidefinite. Chernoff inversion yields
$D-\widehat D\le\sqrt{2x(LD+Q)}+cx$ outside an event of probability
$e^{-x}$. Solving the quadratic inequality gives the stated bound.
\end{proof}

For equal site sizes $n$, the pooled patient mean equals the source-average
mean. With $N=Kn$, define its pooled empirical variance divided by $N$ as
\[
 \hat v_{\rm pool}=\frac{
   \sum_j n(n-1)\hat v_j+\sum_jn(\bar R_j-\bar R_\cdot)^2}{N(N-1)}.
\]
If all patient scores are in $[-M,M]$, Theorem 11 of
\citet{maurer2009}, which permits independent nonidentical observations,
gives the two-sided radius
\[
 u_M=\sqrt{2\hat v_{\rm pool}\log(4/\alpha_M)}
       +\frac{14M\log(4/\alpha_M)}{3(N-1)}.
\]
Thus $J=[\bar R_\cdot\pm\{u_M+B_1+\sqrt{(K-g)U_D/g}\}]$.
Only the pooled-mean implementation requires equal site sizes. The quadratic
bound allows heterogeneous sizes. When $g=K$, no dispersion protection is
needed. These scores are bounded conditional on training; a bound that itself
grows rapidly with fitted weights can still be practically unhelpful.

\paragraph{The exact Gaussian comparator is a different statistic.}
For Gaussian means with common \emph{known} variance $v$, define
$D^\circ=K^{-1}\sum_j(\bar R_j-\bar R_\cdot)^2-(K-1)v/K$.
Then $KD^\circ/v+K-1\sim\chi^2_{K-1}(KD/v)$ exactly.
This identity generally fails with the random subtraction in $\widehat D$.
The implementation reserves exact noncentral chi-square inversion for known
Gaussian variances; bounded scores use the corrected patient statistic.

\subsection{Fourier constraints and all required allowances}

Let $[v_{Lj},v_{Uj}]$ be pilot-certified mean variance bounds, with joint
failure allowance $\delta_\Sigma$ over sites and arms. Known variances are
an oracle option. With $m=\lceil\log(eK)\rceil$, the experimental frequencies are
\[
 t_\ell=\frac{\sqrt\ell}{2\sqrt{\max_jv_{Uj}}},\quad
 \ell=1,\ldots,m,\qquad
 \widehat F(t)=K^{-1}\sum_j e^{t^2v_{Lj}/2}e^{it\bar R_j}.
\]
They are fixed conditional on training and pilot, not selected using the
held-out means. Write $a_j=e^{t^2v_{Lj}/2}$ and
$z=\log(4m/\alpha_F)$. For bounded scores let
\[
 V_F=\sum_ja_j^2\min(1,t^2v_{Uj}),\quad b_F=2\max_ja_j,
 \quad q_F=b_Fz/3.
\]
A stochastic allowance is the smaller of
$2\sqrt{z\sum_ja_j^2}/K$ and
$\sqrt2\{q_F+\sqrt{q_F^2+2V_Fz}\}/K$.
This deterministic choice between two concentration bounds does not spend a
second error allowance. Include the variance-attenuation and non-Gaussian
allowances
\begin{align*}
 e_V(t)&=K^{-1}\sum_j[1-e^{-t^2(v_{Uj}-v_{Lj})/2}],\\
 e_G(t)&=K^{-1}\sum_j\frac{a_j|t|^3}{6n_j^2}
  [R_jn_jv_{Uj}+2\sqrt{2/\pi}(n_jv_{Uj})^{3/2}].
\end{align*}
The latter is a Lindeberg replacement bound using bounded centered third
moments. It is unnecessary for exactly Gaussian means, but cannot be omitted
for causal residual means merely because each has a CLT.

For uniform nuisance drift use $e_B(t)=\rho\min(2,|t|b)$.
For \eqref{eq:h-average}, the valid-subset Taylor expansion yields
\[
 e_B(t)=\min\{2\rho,\rho|t|B_1+t^2B_2/2\}.
\]
In either case, retain $u$ exactly when, for every selected frequency,
\[
 |\widehat F(t)-\rho e^{itu}|
 \le 1-\rho+e_{\rm stochastic}(t)+e_V(t)+e_G(t)+e_B(t).
\]
This defines $F$. Each complex disk is inverted as periodic arcs and their
intersection is computed without a discretized search grid. The hybrid is
localized to $J$; Fourier-only uses the declared residual range $[-1,1]$
in the binary-outcome diagnostic. In a different model a justified domain
must be supplied. This multiscale construction is motivated by Gaussian
contamination inference \citep{wang2026}; the present frequency rule alone
does not prove a sharp length rate in the causal model.

\section{A constructive aggregate nuisance certificate in four strata}

This section concerns the saturated finite-stratum diagnostic, not sparse
high-dimensional calibration. For fixed training fits, let
$\delta_a(x)=m_t^a(x)-m^a(x)$. The nuisance component, including a
counterfactual definition for structurally invalid sites, is
\[
 \gamma_{ja}=\sum_x\{q_j^a(x)p_j(x,A=a)-p_t(x)\}\delta_a(x).
\]
For a structurally valid site $\beta_{ja}=\gamma_{ja}$.
A simultaneous target outcome box bounds $\delta_a$. For each corner $u$,
source pilot patients provide the scores
$Z_{ji}(u)=\ind(A_{ji}=a)q_j^a(X_{ji})u(X_{ji})$.
Use the preceding pooled-mean and exponential-dispersion bounds for their
means $z_j(u)$. Simultaneous target joint $(X,A)$ and marginal $X$
probability intervals, intersected with the simplex and marginal constraints,
bound $\E_tu(X)$ by linear optimization.

If $S(u)$ bounds $|\bar z(u)-\E_tu|$ and $U(u)$ bounds the dispersion of
$z_j(u)$, then every subset $V$ of size $g$ satisfies
\[
 |g^{-1}\sum_{j\in V}\gamma_j(u)|
 \le S(u)+\sqrt{(K-g)U(u)/g},\qquad
 K^{-1}\sum_{j\in V}\gamma_j(u)^2\le U(u)+S(u)^2.
\]
Maximize each bound over the outcome-box corners to obtain $B_1,B_2$.
The absolute linear and squared-norm objectives are convex in $u$, so
corner bounds suffice for the entire box. This protects the unknown valid
subset instead of replacing a maximum by an unprotected average.
The target drift is optimized over the same probability region using
coefficients $[q_t^a\ind(A=a)-1]u(X)$.

With $L=4$ strata there are 16 corners per arm. Allocate $.001$ to the
outcome box, $.001$ to target design probabilities, and $.003$ to pilot
mean/dispersion bounds across the 32 arm-corners. Their sum is
$\delta_B=.005$. Source pilot bounds condition on training; target design
bounds are simultaneous over coefficients and may use all target $(X,A)$
observations. The final union bound does not require these certificate
events to be independent of the target evaluation region.
This computation is linear in $K$ for fixed $L$, but exponential in $L$;
it is not a high-dimensional implementation.

\section{Target region, projection, and explicit coverage}

The target evaluates
$G=(m^1-m^0,\ind(A=1)q_t^1(Y-m^1),\ind(A=0)q_t^0(Y-m^0))^T$.
Let $\hat v_t$ be its mean and $\hat V_t$ its unbiased mean covariance.
Use directions $e_1,e_2,e_3$ and $\ell=(1,1,-1)^T$.
For direction $h$, with deterministic score range $R_h$ and certified
coordinate drift $b_t$, the empirical Bernstein band has radius
\[
 u_h=\sqrt{2h^T\hat V_th\log(16/\alpha_t)}
       +\frac{7R_h\log(16/\alpha_t)}{3(n_t-1)}+|h|^Tb_t.
\]
Their intersection is $C_t$. The TATE band uses the full covariance;
iid patients do not imply independent arms or independent scores on a
shared patient. Known-covariance Gaussian ellipsoids and known-variance
Bernstein regions remain labeled oracle comparisons.

Form
\[
 \mathcal F=\{(d,r_1,r_0)\in C_t:r_a\in C_a\},\qquad
 I=\operatorname{hull}\{d+r_1-r_0:(d,r_1,r_0)\in\mathcal F\}.
\]
If the intersection is empty, use the prescribed shorter fallback: form
$J_S=[\hat d\pm u_d]+J_1-J_0$, tighten using the source-set hulls when
both are nonempty, and choose the shorter of that source interval and the
prespecified target interval. The reported length never exceeds $|J_S|$.
For Fourier-only, substitute the declared residual domain for $J_a$.
A final projection onto $[-1,1]$ is allowed for binary-outcome TATE.
All fallback repetitions are retained. The interval midpoint and a projected
source-center point are secondary summaries, not optimally weighted estimates.

\begin{proposition}[Hybrid coverage reduction]
If the stated score independence/range conditions, valid-count assumption,
and nuisance and variance certificates hold with their allocated probabilities,
and
\[
 \alpha_t+\delta_B+\delta_\Sigma+
 \sum_{a=0}^1(\alpha_{F,a}+\alpha_{M,a}+\alpha_{D,a})\le\alpha,
\]
then the reported interval covers $\tau_t$ with probability at least $1-\alpha$.
An omitted module receives zero error or an explicitly reserved unused allowance.
\end{proposition}
\begin{proof}
On the intersection of all certificate and concentration events,
$r_a\in J_a\cap F_a$ and $(d,r_1,r_0)\in C_t$. Hence the feasible set is
nonempty and its projection contains $\tau_t$. The union bound gives the
result without independence between arms or modules. A fallback is used only
outside this simultaneous good event and cannot invalidate this argument.
\end{proof}

The paired fitted experiment uses $.02$ for the target, $.005$ each for
variance and nuisance certification, and $.01$ per source arm. The hybrid
splits each arm allowance equally between Fourier and dispersion; dispersion
splits its allowance between mean and energy. Single-component ablations use
the whole arm allowance. The direct-dispersion ablation still reserves the
unused $.005$ variance allowance, so the comparison does not benefit from
reassigning it after observing results.

\section{What this does and does not establish for growing source counts}

The finite-sample statements allow arbitrary local structural biases and
any finite $K$. Thus they yield coverage along growing-$K$ sequences when
the stated ranges, independence, valid-count, and certificate assumptions
hold along that sequence. This is not yet a high-dimensional fitted RoCE
coverage theorem or an optimality claim.

With comparable $n_j=n$, uniformly bounded ranges, equal conditional source means ($D=0$), and
fixed error allocations, $L=O(1/(nK))$, $Q=O(1/(n^2K))$, and
$c=O(1/(n^2K))$. Consequently $U_D=O_p(1/(n\sqrt K))$ and the dispersion
protection is $O_p(n^{-1/2}K^{-1/4})$ when $g/K$ stays bounded away from zero.
If $g$ is fixed this conclusion changes. Nonvanishing contamination dispersion,
shared target prediction noise, and nuisance budgets can prevent useful
shrinkage. Fourier's additional length improvement requires a separate
argument excluding wrong centers at the required scales; a coverage proof
alone does not supply it. The Gaussian adaptive-interval results of
\citet{wang2026} are benchmarks, not transferred lower bounds.

The primary next decision is empirical and theoretical: do direct calibration
moments and dispersion bounds make borrowing improve on an equally improved
target-only procedure? Preserve all earlier experiments and compare both
ordinary target-normal and finite-sample target references. Gains over a
looser historical reference are not sufficient evidence of source efficiency.
